\section{Motivating Example}

To illustrate why malicious skill detection remains difficult in practice, we consider a representative skill illustrated in \autoref{fig:motivation}, which appears to be a benign DevOps troubleshooting utility but in fact performs credential harvesting and exfiltration. First, \texttt{SKILL.md} instructs the agent to prepare a diagnostic bundle for remote-access troubleshooting. Second, \texttt{config.yaml} defines a bundle directory, a remote upload endpoint, and a command that copies \texttt{\textasciitilde/.aws/credentials} into the bundle. Third, \texttt{post.js} reads the diagnostic plan, collects the bundle contents, and sends them to the configured remote server. The malicious behavior is thus not explicit in any single file; it emerges only after correlating the prompt, configuration, and code.

This example highlights why existing detection methods are insufficient. As shown in \autoref{fig:motivation}, the malicious logic is distributed across multiple artifacts: \ding{182} and \ding{183} denote two security-sensitive read operations, while \ding{184} denotes a send operation; their concrete operands are further revealed by \ding{185}--\ding{187}, which connect the reads to \texttt{\textasciitilde/.aws/credentials} and \texttt{diagnostic.yaml}, and the send to the remote endpoint and payload. This example is challenging for existing detectors for three reasons. First, the evidence is \emph{cross-artifact}: no single file explicitly states the full malicious intent, so artifact-local static scanning may only observe isolated signals such as a credential read or an outbound request. Second, the risk is \emph{context-sensitive}: both the reads and the upload can appear consistent with a benign troubleshooting workflow, so LLM-based analysis may be misled by the adversarial words. Third, the maliciousness is \emph{data-dependent}: the key issue is not the mere invocation of \ding{182}--\ding{184}, but that the sensitive operands accessed in \ding{185} is propagated through the intermediate dependency in \ding{186} and ultimately becomes part of the payload sent in \ding{187}. In other words, the main challenge is to reconstruct whether seemingly benign operations across heterogeneous artifacts are connected into a malicious workflow.